\SourceLeaf{290}%
\[
\begin{aligned}
\mathrm{d}\mathord{.}\frac{\operatorname{tang}\psi}{m}&=-\frac{l\mathrm{d}x}{m}\\
\mathrm{d}u&=\operatorname{tang}\psi\mathord{.}\mathrm{d}x
\end{aligned}
\]
so ergibt sich
\[
\begin{aligned}
\operatorname{tang}\psi&=\mathfrak{A}(1+\mu x^3+\mu'x^4+\text{usw.})-\tfrac{1}{3}\lambda x^3-\tfrac{1}{4}\lambda'x^4-\text{usw.}\\
u&=\mathfrak{B}+\mathfrak{A}(x+\tfrac{1}{4}\mu x^4+\tfrac{1}{5}\mu'x^5+\text{usw.})-\tfrac{1}{12}\lambda x^4-\tfrac{1}{20}\lambda'x^5-\text{usw.}
\end{aligned}
\]
\indent Die durch die Integration eingeführten Konstanten \(\mathfrak{A},\mathfrak{B}\) lassen sich aus der Bedingung bestimmen, dass \(u=0\) für die beiden Werte von \(x\) sein muss, die den Punkten \(F,G\) entsprechen. Es seien dies \(x=-\tfrac{1}{2}(h-\delta)\) und \(x=+\tfrac{1}{2}(h+\delta)\). Dabei bezeichnet \(\delta\) den Wert von \(2x\) in der Mitte zwischen \(F\) und \(G\) und ist im Allgemeinen eine Größe derselben Ordnung wie \(h\), oder höherer Ordnung, wenn \(H\) sehr nahe an dieser Mitte liegt. Daraus lässt sich leicht der folgende, bis einschließlich zur Ordnung \(h^3\) genaue Ausdruck für \(\mathfrak{A}\) herleiten
\[
\mathfrak{A}=\frac{\lambda((h+\delta)^4-(h-\delta)^4)}{192h}=\tfrac{1}{24}\lambda\delta(hh+\delta\delta)
\]
\indent Man setzt diesen Ausdruck in die Reihe für \(\operatorname{tang}\psi\) ein. Gibt man dann der Veränderlichen \(x\) die bestimmten Werte \(-\tfrac{1}{2}(h-\delta),+\tfrac{1}{2}(h+\delta)\), so ergibt sich, ebenfalls bis zur dritten Ordnung genau,
\[
\begin{aligned}
\operatorname{tang}\psi^0&=\tfrac{1}{24}\lambda h(hh-2h\delta+3\delta\delta)\\
\operatorname{tang}\psi'&=-\tfrac{1}{24}\lambda h(hh-2h\delta+3\delta\delta)
\end{aligned}
\]
\indent Im besonderen Fall der später auszuarbeitenden Untersuchung kommt zu der oben genannten Bedingung noch hinzu, dass der Normalparallelkreis genau in der Mitte zwischen den beiden Parallelkreisen liegt, auf denen sich die Punkte \(F,G\) befinden. Daher sind schon die gekürzten Ausdrücke
\[
\begin{aligned}
\operatorname{tang}\psi^0&=\tfrac{1}{24}\lambda h^3\\
\operatorname{tang}\psi'&=-\tfrac{1}{24}\lambda h^3
\end{aligned}
\]
bis zur dritten Ordnung genau, wie sich leicht auf folgende Weise zeigen lässt. Bezeichnet man die Breite von \(F\) mit \(Q+q\) und die von \(G\) mit \(Q-q\), so ergeben sich aus den sphärischen Dreiecken (\(F,H\), Pol) und (\(G,H\), Pol) die Gleichungen
